\section{\add{Previous frequency indexing convention}\label{sec:apx:fmoc1}}

\add{This section describes the frequency axis indexing scheme that was initially explored before
the current standardized version. Documenting it here serves two purposes:
\begin{itemize}
    \item to document earlier versions of FMOC and SFMOC files that may still exist;
    \item to record an already explored solution and thus avoid attempting to reintroduce it in the future.
\end{itemize}}

\add{The original idea was to (almost) directly use, as the index at the deepest level,
the integer representation of the 64-bit floating-point value of a frequency expressed in Hz.
Performing MOC operations (except \textit{degrade}) on floating-point ranges is equivalent to
performing them on the corresponding integer representation ranges because:
\begin{itemize}
  \item frequencies are strictly positive values;
  \item for positive floating-point numbers, the numerical ordering is identical to the ordering of their
        integer bit representations.
\end{itemize}}

\add{The only difference was that the exponent range of a 64-bit floating-point value was considered too
large for the intended use. Therefore, it was restricted to raw exponent values
$\in [929, 1184]$ instead of the full IEEE~754 range $\in [0, 2047]$.
The selected interval contains 255 distinct values and therefore requires 8 bits
instead of the 11 bits used in standard double-precision representation.
Using the 52 mantissa bits together with these 8 exponent bits leads to 60 hierarchical
orders, with a maximum order of 59. This also defines the, somewhat arbitrary, validity range
$f \in [5.0487097934144756 \times 10^{-29},\; 5.846006549323611 \times 10^{48}[$.
The following pseudo-code describes the transformation, which replaces the
11-bit exponent $\in [929,1184]$ of a floating-point value with an 8-bit exponent
$\in [0,255]$.}
\lstdefinelanguage{Pseudocode}{
  keywords={input, output, return, datatype, function, in, if, else, foreach, while, begin, end
  },
  sensitive=false,
  comment=[l]{//},
  morecomment=[s]{/*}{*/},
  morestring=[b]',
  morestring=[b]"
}
\begin{lstlisting}[language=Pseudocode,mathescape=true,caption=Frequency in Hz to order 59 index routine,label={code:freq2index}]
Input: long float freq_hz
Output: long int index
begin
  // Set the 64-bits float exponent bits mask
  long int emask $\gets$ 0x7FF << 52
  // Set the mantissa (+sign) bits mask
  long int mmask $\gets$ !emask
  // Get the bits level representation of the input frequency
  long int index = freq_hz.to_bits();
  // Extract the exponent from the input frequency
  long int exp $\gets$ (index & emask) >> 52;
  // Replace the exponent value in [929, 1184] by a value in [0, 255]  
  exp $\gets$ exp - 929
  // Replace the 11 exponent bits by the new 8 exponent bits
  index $\gets$ (index & mmask) | (exp << 52)
  return index
end
\end{lstlisting}


\begin{lstlisting}[language=Pseudocode,mathescape=true,caption=Frequency in Hz from order 59 index routine,label={code:index2freq}]
Input: long int index
Output: long float freq_hz
begin
  // Set the 64-bits float exponent bits mask
  long int emask $\gets$ 0x7FF << 52
  // Set the mantissa (+sign) bits mask 
  long int mmask $\gets$ !emask
  // Extract the exponent part in [0, 255]
  long int exp = (index & emask) >> 52;
  // Replace epxonent value in [0, 255] by value in [929, 1184]
  exp $\gets$ exp + 929
  // Replace 8 exponent bits by 11 exponent bits
  index $\gets$ (index & mmask) | (exp << 52)
  // Transform bit representation into float value
  return index.f64_value_from_bits()
end
\end{lstlisting}


\add{One advantage of this approach was the existence of a perfect bijection between the
64-bit floating-point representation of a frequency (in Hz) within the valid range and
its index at the deepest resolution. The conversions were also extremely fast, requiring
only a few bitwise operations. In addition, the number of orders (60) was not limited to
the number of mantissa bits.}

\add{However, a major issue arose from the representation of the axis using WCS, as used in HiPS 2.0.
Because the representation allocates 8 bits to the exponent and 52 bits to the mantissa,
the resulting WCS cannot be expressed purely as either linear or logarithmic.
For orders 0 to 7, the WCS behaves logarithmically.
For orders 8 to 59, inside a cell or order at least 7, it behaves linearly.
For ranges spanning order 7, however, the WCS becomes a mixture of logarithmic and linear
components, which cannot be represented easily.
This limitation was the primary motivation for abandoning this approach.}





